10.4 Graphing Sine and Cosine Functions

1. Key Info from Equation  y = a sin b(x-h) + k
Amplitude = |a|, reflects if a < 0
$$ \text{Period} = \frac{2\pi}{|b|} \qquad \text{Midline: } y = k $$
h = horizontal (phase) shift, k = vertical shift

2. Quick ID Examples
$$ g(x) = 4\sin x: \quad |a|=4,\ T=2\pi,\ h=0,\ k=0 $$
$$ g(x) = \tfrac{1}{2}\cos 2\pi x: \quad |a|=\tfrac{1}{2},\ T=1,\ h=0,\ k=0 $$
$$ g(x) = 2\sin 4x+3: \quad |a|=2,\ T=\tfrac{\pi}{2},\ k=3 $$
$$ g(x) = 5\cos\tfrac{1}{2}(x-3\pi): \quad |a|=5,\ T=4\pi,\ h=3\pi $$
$$ g(x) = -2\sin\tfrac{2}{3}(x-\tfrac{\pi}{2}): \quad |a|=2,\ T=3\pi,\ h=\tfrac{\pi}{2},\ \text{reflected} $$

3. Five Key Points
Divide one period into 4 equal parts starting at h.
sin: starts (0,0) max min pattern
cos: starts (0,a) down up pattern
Shift all x by h, all y by k. If a < 0, flip max/min.

4. Reading a Graph (MC tips)
Amplitude = (max - min) / 2
Period = x-distance for one full cycle
Midline k = (max + min) / 2
Phase shift: where the pattern starts vs x=0
Negative a: sin starts going DOWN, cos starts at BOTTOM
